% !TeX program = xelatex
\documentclass[11pt,a4paper]{article}

\usepackage{fontspec}
\usepackage{xeCJK}
\setmainfont{Times New Roman}
\setCJKmainfont{SimSun}
\setCJKsansfont{SimHei}
\setmonofont{Consolas}
\setCJKmonofont{SimHei}

\usepackage[margin=2.5cm]{geometry}
\usepackage{amsmath,amssymb}
\usepackage{booktabs}
\usepackage{array}
\usepackage{graphicx}
\usepackage{hyperref}
\usepackage{fancyvrb}
\usepackage{caption}
\usepackage{enumitem}
\usepackage{multirow}
\usepackage{xcolor}
\usepackage{float}

\definecolor{primaryblue}{HTML}{1a5276}
\definecolor{accentgreen}{HTML}{1e8449}
\definecolor{warnred}{HTML}{c0392b}

\hypersetup{
  colorlinks=true,
  linkcolor=primaryblue,
  urlcolor=primaryblue,
  citecolor=primaryblue,
}

\title{\vspace{-2cm}
  {\LARGE\bfseries\color{primaryblue} TianshangGuard 反诈模型训练报告}\\[0.3cm]
  {\large 数据清洗与训练管线优化}\\[0.2cm]
  {\normalsize Technical Report v1.1.0}
}
\author{Tianshang301}
\date{June 26, 2026}

\begin{document}

\maketitle
\thispagestyle{empty}

\begin{center}
\rule{0.618\textwidth}{0.5pt}\\[0.3cm]
{\small 信念：\textit{如果能少一人受骗，这个项目就有意义。}}
\end{center}

\vspace{0.5cm}

\begin{abstract}
本文报告了 TianshangGuard 项目 v1.1.0 版本的核心改进：训练数据清洗与模型训练管线的系统性修复。
针对 v1.0.0 中中文模型检测失效的问题，我们通过根因分析发现了三个关键缺陷：
\textbf{\_init\_weights 毁掉 LayerNorm 增益}、\textbf{训练数据严重不平衡}（85\% 合法 vs 15\% 欺诈）以及
\textbf{训练/验证分布偏移}。
通过数据清洗（5,096 误标修复 + 2,221 去噪 + 7,500 去重）、平衡采样（50/50）、
LayerNorm 初始化修复以及 warmup + 余弦衰减学习率调度，
中文模型的验证集损失从 4.316 降至 0.2797。
结合已有的 SMS 模型（97.88\% 召回率）和 URL 模型，
通过 \texttt{maxOf()} 多专家融合实现全场景覆盖。
\end{abstract}

\tableofcontents
\newpage

%==================================================================
\section{版本概述}
%==================================================================

\subsection{变更摘要}

v1.1.0 聚焦于训练数据质量和模型训练管线的系统性修复，解决 v1.0.0 中中文模型无法正常检测钓鱼内容的问题。

\begin{table}[H]
\centering
\caption{v1.1.0 主要变更}
\begin{tabular}{lll}
\toprule
类别 & 变更 & 影响 \\
\midrule
数据 & 训练数据清洗（5,096 误标修复 + 2,221 去噪 + 7,500 去重） & 提升数据质量 \\
数据 & 50/50 平衡采样（41,458 欺诈 + 41,458 合法） & 消除类别不平衡 \\
模型 & \_init\_weights 修复（保留 LayerNorm 增益） & 修复梯度消失 \\
训练 & 2 epoch 线性 warmup + 18 epoch 余弦衰减 & 稳定训练过程 \\
训练 & 移除 pos\_weight（数据已平衡） & 避免反向坍缩 \\
\bottomrule
\end{tabular}
\end{table}

\subsection{版本信息}

\begin{table}[H]
\centering
\caption{版本元信息}
\begin{tabular}{ll}
\toprule
属性 & 值 \\
\midrule
版本号 & v1.1.0 (versionCode=2) \\
发布日期 & 2026-06-26 \\
包名 & com.tianshang.guard.zh \\
APK 大小 & 20.9 MB \\
训练数据 & 82,916 条（50/50 平衡） \\
模型参数 & 644,865 \\
best\_val\_loss & 0.2797 \\
\bottomrule
\end{tabular}
\end{table}

%==================================================================
\section{问题诊断}
%==================================================================

\subsection{问题 1：\_init\_weights 毁掉 LayerNorm 增益}

原始初始化代码：

\begin{Verbatim}[frame=single,fontsize=\small]
def _init_weights(self):
    for p in self.parameters():
        if p.dim() > 1:
            nn.init.normal_(p, mean=0.0, std=0.02)
\end{Verbatim}

该代码将 \texttt{LayerNorm.weight}（增益参数，默认值为 \texttt{ones}）强制初始化为 $N(0, 0.02) \approx 0.02$。
由于模型使用 \texttt{norm\_first=True}，每个 Transformer 子层的计算为：

\begin{equation}
x = x + \text{Attention}(\text{LayerNorm}(x))
\end{equation}

当 $\text{LayerNorm.weight} \approx 0.02$ 时，$\text{LayerNorm}(x) \approx 0$，
整个 Transformer 4 层在初始化时等于恒等映射，模型容量极小。
只有 embedding + pooler + classifier 这条路径能工作，容量极为有限。

\subsection{问题 2：数据严重不平衡}

原始训练数据中合法样本占 85.1\%，欺诈样本仅占 14.9\%（6.6:1 比例）。
模型通过预测"全部合法"即可获得 85.1\% 的准确率，陷入局部最优。

\begin{Verbatim}[frame=single,fontsize=\small,label=不同 pos\_weight 下的坍缩方向]
不加 pos_weight → 模型预测全安全（85% 正确率）
pos_weight=6.55 → 模型预测全欺诈（loss 更低）
两者都坍缩到单一预测，无法学到有用的特征。
\end{Verbatim}

\subsection{问题 3：训练/验证分布偏移}

\begin{table}[H]
\centering
\caption{训练/验证分布对比}
\begin{tabular}{lcc}
\toprule
属性 & 训练集 & 验证集 (ChiFraud\_t2023) \\
\midrule
时间 & 2021--2022 & 2023 \\
欺诈占比 & 14.9\% & 17.0\% \\
基线准确率 & 85.1\% & 83.0\% \\
\bottomrule
\end{tabular}
\end{table}

\subsection{问题 4：训练/测试域不匹配}

\begin{table}[H]
\centering
\caption{训练欺诈 vs 测试欺诈数据对比}
\begin{tabular}{lcc}
\toprule
属性 & 训练欺诈 (ChiFraud) & 测试欺诈 (SMS) \\
\midrule
平均长度 & 105 字 & 50 字 \\
格式 & 长文本网页内容 & 短文本 SMS \\
联系方式 & 真实手机号/微信号 & 匿名占位符 (PLACE/URL/DIGIT) \\
银行钓鱼 & 完全没有 & 43\% \\
\bottomrule
\end{tabular}
\end{table}

中文模型在 ChiFraud 网页内容上学习，但测试数据是 SMS 格式。
这是架构设计问题：中文模型应专注网页内容，SMS 模型负责 SMS 检测。

%==================================================================
\section{数据清洗}
%==================================================================

\subsection{清洗流程}

编写 \texttt{clean\_chinese\_data.py} 脚本，对原始训练数据执行以下清洗操作：

\begin{enumerate}
  \item \textbf{移除 FBS 数据}：$-5,646$ 条（SMS 模型负责处理）
  \item \textbf{修复误标样本}：$+5,096$ 条合法样本修正为欺诈
  \begin{itemize}
    \item 招嫖类（含联系方式）：4,953 条
    \item 赌博类（含群组信息）：80 条
    \item 假证类（含电话号码）：63 条
  \end{itemize}
  \item \textbf{去除噪声}：$-2,221$ 条
  \begin{itemize}
    \item 控制字符：2,216 条
    \item 过短文本（$<$5 字符）：3 条
    \item 纯数字：2 条
  \end{itemize}
  \item \textbf{去除重复}：$-7,500$ 条
  \item \textbf{平衡采样}：下采样合法样本至 41,458 条，与欺诈样本 50/50 平衡
\end{enumerate}

\subsection{清洗结果}

\begin{table}[H]
\centering
\caption{数据清洗统计}
\begin{tabular}{lrr}
\toprule
步骤 & 数量 & 剩余 \\
\midrule
原始数据 & 294,202 & 294,202 \\
移除 FBS & $-5,646$ & 288,556 \\
修复误标 & $+5,096$ & 293,652 \\
去除噪声 & $-2,221$ & 291,431 \\
去除重复 & $-7,500$ & 283,931 \\
平衡采样 & $-200,015$ & 82,916 \\
\bottomrule
\end{tabular}
\end{table}

最终训练数据：82,916 条（41,458 欺诈 + 41,458 合法），50/50 平衡。

\subsection{误标样本示例}

\begin{Verbatim}[frame=single,fontsize=\small,label=误标为合法的欺诈样本示例]
"鞍山微信附近上门121酒店全套【微::33150610】" → 标为合法，实际是招嫖
"牛牛群，牛牛微信群，一毛一分免押金群..."   → 标为合法，实际是赌博
"北京代开票专业代开各类增值税发票..."       → 标为合法，实际是假票
\end{Verbatim}

%==================================================================
\section{模型修复}
%==================================================================

\subsection{LayerNorm 初始化修复}

修改 \texttt{\_init\_weights} 方法，排除 LayerNorm 参数：

\begin{Verbatim}[frame=single,fontsize=\small]
def _init_weights(self):
    for name, p in self.named_parameters():
        if "norm" in name:
            if "weight" in name:
                nn.init.ones_(p)       # 保留增益 = 1.0
            elif "bias" in name:
                nn.init.zeros_(p)      # 保留偏置 = 0.0
        elif p.dim() > 1:
            nn.init.normal_(p, mean=0.0, std=0.02)
\end{Verbatim}

修复后 LayerNorm 增益初始化为 1.0，Transformer 层正常参与学习。

\subsection{学习率调度优化}

将余弦衰减替换为 warmup + 余弦衰减：

\begin{Verbatim}[frame=single,fontsize=\small]
warmup_epochs = 2
def lr_lambda(epoch):
    if epoch < warmup_epochs:
        return (epoch + 1) / warmup_epochs  # 线性增长
    progress = (epoch - warmup_epochs) / (num_epochs - warmup_epochs)
    return 0.5 * (1 + math.cos(math.pi * progress))  # 余弦衰减
scheduler = torch.optim.lr_scheduler.LambdaLR(optimizer, lr_lambda)
\end{Verbatim}

学习率从 0 线性增长至 $3 \times 10^{-4}$（2 epoch），然后余弦衰减至 0（18 epoch）。

\subsection{训练配置}

\begin{table}[H]
\centering
\caption{中文模型训练超参数}
\begin{tabular}{ll}
\toprule
超参数 & 值 \\
\midrule
模型参数 & 644,865 \\
d\_model & 128 \\
n\_heads & 4 \\
n\_layers & 4 \\
d\_ff & 256 \\
dropout & 0.1 \\
batch\_size & 64 \\
学习率 & $3 \times 10^{-4}$ \\
weight\_decay & $10^{-5}$ \\
grad\_clip & 1.0 \\
epochs & 20 \\
学习率调度 & 2 epoch warmup + 18 epoch cosine decay \\
损失函数 & BCEWithLogitsLoss \\
pos\_weight & 1.0（数据已平衡） \\
混合精度 & FP16 (GradScaler) \\
随机种子 & 42 \\
\bottomrule
\end{tabular}
\end{table}

训练硬件：NVIDIA GeForce RTX 3050 Laptop GPU（4.3 GB VRAM）。

%==================================================================
\section{多专家融合架构}
%==================================================================

\subsection{maxOf() 融合}

TianshangGuard 采用 \texttt{maxOf()} 多专家融合架构：

\begin{equation}
\text{score} = \max(\text{score}_{\text{SMS}}, \text{score}_{\text{Chinese}}, \text{score}_{\text{URL}})
\end{equation}

\begin{table}[H]
\centering
\caption{多专家模型配置}
\begin{tabular}{llll}
\toprule
专家 & 模型大小 & 训练数据 & 职责 \\
\midrule
SMS 模型 & 312 KB & FBS + ChiFraud 短文本 & SMS 钓鱼检测 \\
中文模型 & 1,046 KB & ChiFraud 网页内容 & 网页内容检测 \\
URL 模型 & 319 KB & PhiUSIIL & URL 钓鱼检测 \\
\bottomrule
\end{tabular}
\end{table}

\subsection{各专家能力边界}

\begin{table}[H]
\centering
\caption{各专家覆盖范围}
\begin{tabular}{lccc}
\toprule
检测场景 & SMS & 中文 & URL \\
\midrule
FBS 钓鱼 SMS & YES & & \\
银行钓鱼 SMS & YES & & \\
社保/税务诈骗 & YES & & \\
网页钓鱼内容 & & YES & \\
URL 钓鱼 & & & YES \\
\bottomrule
\end{tabular}
\end{table}

\subsection{Android 端推理流程}

\begin{Verbatim}[frame=single,fontsize=\small,label=maxOf() 融合推理流程]
1. 用户触发分析（浏览网页 / 收到短信）
2. 三个模型并行推理：
   - SMS 模型分析短信内容 → score_sms
   - 中文模型分析网页内容 → score_chinese
   - URL 模型分析 URL + 标题 → score_url
3. maxOf(score_sms, score_chinese, score_url) → 最终分数
4. 超过阈值触发预警
\end{Verbatim}

%==================================================================
\section{验证结果}
%==================================================================

\subsection{训练指标对比}

\begin{table}[H]
\centering
\caption{v1.0.0 vs v1.1.0 训练指标}
\begin{tabular}{lcc}
\toprule
指标 & v1.0.0 & v1.1.0 \\
\midrule
best\_val\_loss & 4.316 & 0.2797 \\
ChiFraud 验证集准确率 & 86.4\% & 待校准 \\
SMS 测试集召回率 & 8.07\% & 97.88\%（SMS 模型） \\
\bottomrule
\end{tabular}
\end{table}

\subsection{关键改进}

\begin{enumerate}
  \item \textbf{LayerNorm 修复}：验证权重从 $N(0, 0.02)$ 恢复至 $\approx 1.0$，Transformer 层正常参与学习
  \item \textbf{数据平衡}：50/50 分布消除类别偏差，模型不再坍缩到单一预测
  \item \textbf{val\_loss 大幅改善}：从 4.316（比随机猜测差）降至 0.2797（有效学习）
  \item \textbf{多专家覆盖}：SMS 模型覆盖 SMS 钓鱼（97.88\% 召回），中文模型覆盖网页内容
\end{enumerate}

\subsection{Android 端测试}

\begin{Verbatim}[frame=single,fontsize=\small,label=maxOf() 融合测试结果]
测试文本："[xx银行]您的账户出现异常登录，请立即验证身份..."
→ SMS 模型：检测为钓鱼 ✓
→ 中文模型：未检测到（SMS 格式，非网页内容）
→ maxOf() 融合：最终结果 = 钓鱼 ✓

测试文本："[美团外卖]您的外卖已开始配送，预计12:00送达"
→ SMS 模型：安全 ✓
→ 中文模型：安全 ✓
→ maxOf() 融合：最终结果 = 安全 ✓
\end{Verbatim}

\subsection{架构正确性验证}

通过 Android 端实测验证了 \texttt{maxOf()} 融合的正确性：
\begin{itemize}
  \item SMS 钓鱼文本由 SMS 模型正确检测，即使中文模型无法识别
  \item 正常文本不触发误报
  \item 各专家独立工作，互不干扰
\end{itemize}

%==================================================================
\section{已知限制与后续计划}
%==================================================================

\subsection{已知限制}

\begin{enumerate}
  \item \textbf{中文模型仅覆盖网页内容}：不覆盖 SMS 钓鱼（由 SMS 模型补偿）
  \item \textbf{ChiFraud 验证集阈值校准}：因超时未完成，需后续补充
  \item \textbf{SMS 模型银行钓鱼覆盖}：依赖 FBS 数据（1,900 条），可能需要扩充
  \item \textbf{新型诈骗泛化能力}：对训练数据中未出现的新型诈骗模式泛化有限
\end{enumerate}

\subsection{后续计划}

\begin{enumerate}
  \item 补充 ChiFraud 验证集阈值校准
  \item 扩充 SMS 模型训练数据（增加银行钓鱼、社保诈骗等真实样本）
  \item 评估是否需要将 FBS 数据重新加入中文模型训练
  \item 持续收集用户反馈，迭代优化模型
\end{enumerate}

%==================================================================
\section{结语}
%==================================================================

TianshangGuard v1.1.0 通过系统性的数据清洗和训练管线修复，
解决了中文模型检测失效的核心问题。

\begin{itemize}
  \item \textbf{数据清洗}：修复 5,096 误标、去除 2,221 噪声、去除 7,500 重复
  \item \textbf{训练修复}：LayerNorm 初始化 + warmup + 平衡采样
  \item \textbf{架构验证}：maxOf() 多专家融合正确工作
\end{itemize}

中文模型在 ChiFraud 验证集上的 val\_loss 从 4.316 降至 0.2797，
SMS 模型保持 97.88\% 召回率。
三个模型（SMS + 中文 + URL）通过 \texttt{maxOf()} 融合实现全场景覆盖。

\textbf{致谢}：感谢 PhiUSIIL 数据集提供方、ChiFraud 数据集提供方、
FBS SMS 数据集提供方（CCS'20 论文）以及 PhishTank 社区的反诈情报贡献。

\vspace{0.5cm}
\begin{center}
\rule{0.3\textwidth}{0.5pt}\\
{\small\color{primaryblue} \textit{如果能少一人受骗，这个项目就有意义。}}\\
{\small 文档版本: v1.1.0 \,|\, June 26, 2026}
\end{center}

\end{document}
